\section{Fixed-cofactor parabolas}
\label{sec:cofactor-curves}
%======================================================================

Once the cofactor is fixed, the canonical points are no longer an unstructured cloud.  They lie on explicit curves, which lets arithmetic transitions be read geometrically.

Fix a positive cofactor $c$ and let $q$ vary.  In squared Cartesian coordinates,
\begin{equation}
 x=cq,
 \qquad
 y=\frac{q^2-c^2}{2}.
 \label{eq:curve-param}
\end{equation}
Eliminating $q=x/c$ gives the cofactor parabola
\begin{equation}
 \boxed{
 \gamma_c(x)
 =\frac{x^2-c^4}{2c^2}.
 }
 \label{eq:gamma-c}
\end{equation}
The canonical point cloud is not the filled region between envelopes.  It is the arithmetic sampling
\begin{equation}
 x=cq,
 \qquad
 q\text{ prime},
 \qquad
 q>\Pplus(c),
 \label{eq:curve-sampling}
\end{equation}
on the family \cref{eq:gamma-c}.

The curve crosses the real axis at
\begin{equation}
 x=c^2.
 \label{eq:curve-axis-cross}
\end{equation}
Points to the right have $q>c$ and positive ordinate; points to the left have $q<c$ and negative ordinate.

The $c=1$ curve is
\begin{equation}
 \gamma_1(x)=\frac{x^2-1}{2}.
 \label{eq:prime-curve}
\end{equation}
Its arithmetic samples are exactly the prime points, since $m=1\cdot q=q$.  It is the outer positive envelope of the complete canonical cloud.  The curve
\begin{equation}
 \gamma_3(x)=\frac{x^2-81}{18}
 \label{eq:three-curve}
\end{equation}
is the first nontrivial odd-cofactor envelope used in semiprime search.

\begin{figure}[htbp]
\centering
\begin{tikzpicture}
\begin{axis}[
  width=0.92\textwidth,
  height=0.55\textwidth,
  xmin=0,xmax=60,
  ymin=-20,ymax=1850,
  xlabel={$x=\Repart w=cq$},
  ylabel={$y=\Impart w=(q^2-c^2)/2$},
  grid=major,
  legend style={at={(0.98,0.98)},anchor=north east,fill=white,draw=gray!40},
  samples=240,
  domain=0:60,
  clip=true
]
  \addplot[very thick,rust] {(x^2-1)/2};
  \addlegendentry{$\gamma_1$}
  \addplot[very thick,navy] {(x^2-81)/18};
  \addlegendentry{$\gamma_3$}
  \addplot[very thick,teal] {(x^2-625)/50};
  \addlegendentry{$\gamma_5$}
  \addplot[very thick,black,dashed] {(10000-x^2)/200};
  \addlegendentry{$\Gamma_{10}$}
  \addplot[only marks,mark=*,mark size=2.2pt,black] coordinates {(10,49.5) (30,45.5) (50,37.5)};
  \node[anchor=south west,font=\small] at (axis cs:10,49.5) {$c=1$};
  \node[anchor=south west,font=\small] at (axis cs:30,45.5) {$c=3$};
  \node[anchor=south west,font=\small] at (axis cs:50,37.5) {$c=5$};
\end{axis}
\end{tikzpicture}
\caption{Cofactor parabolas $\gamma_c$ and the moving cutoff parabola $\Gamma_R$ for $R=10$.  Each $\gamma_c$ meets $\Gamma_R$ at $x=cR$.  Actual canonical points are discrete prime samples on the curves.}
\label{fig:cofactor-parabolas}
\end{figure}

%======================================================================
